%%%%%%%%%%%%%%%%%%%%%%%%%%%%%%%%%%%%%%%%%%%%%%%%%%%%%%%%%%%%
\section{Related Work}
\label{sec:related-work}

Recent surveys~\cite{acm-ai-agents-survey2024,agentic-security-survey2025,gong-prompt-injection2024,liu-protocol-exploits2025} identify four critical gaps in agentic AI security: unpredictable multi-step workflows, cross-framework compositions, variable environments, and untrusted entity interactions. These surveys diagnose fragmentation but lack systems solutions. Prior work addresses either application semantics or infrastructure primitives; we provide a trust layer positioned between both—protocol-level primitives protecting four fundamental control surfaces (prompts, tools, data, context) where all agent-resource interactions occur.

\textbf{Application-Layer Defenses Fail at Composition.} Framework-specific security operates in isolation. LangChain provides schemas and validation~\cite{langchain2023}; MCP implements authentication and permissions~\cite{mcp2024}; OpenAI enforces rate limits and moderation~\cite{openai-assistants2024}. Recent vulnerabilities demonstrate composition gaps: CVE-2024-8309 (LangChain prompt injection→SQL, CVSS 9.0)~\cite{cve-2024-8309}, CVE-2024-36480 (RCE via unsanitized eval())~\cite{cve-2024-36480}, CVE-2025-6514 (MCP supply chain, 558K+ downloads)~\cite{github-mcp-cve}, Atlas browser attack (October 2025)~\cite{atlas-cve,openai-ciso-atlas}. No cryptographic binding exists across framework boundaries.

Recent academic work addresses prompt injection defenses~\cite{perez2022ignore,liu2023prompt,jain2023baseline}, tool use safety~\cite{ruan2023identifying}, and RAG retrieval poisoning~\cite{zou2024poisonedrag,cheng2024parden}—malicious instructions embedded in documents that compromise reasoning. Semantic defenses (guardrails~\cite{nemoguardrails2023,rebuff2023,llamaguard2023}) achieve 60-80\% detection but remain bypassable~\cite{wei2023jailbroken,zou2023universal}. We provide deterministic guarantees (100\% recall, zero false positives) through defense in depth: by-design elimination renders attacks cryptographically infeasible; by-policy enforcement blocks violations. The approaches are complementary—semantic defenses reduce attack surface; cryptographic enforcement provides non-bypassable verification.

\textbf{Infrastructure Primitives Lack Workflow Semantics.} Infrastructure primitives provide strong guarantees but lack workflow semantics. Identity systems (SPIFFE~\cite{spiffe2023}, AWS IAM~\cite{aws-iam2023}) and policy engines (OPA~\cite{opa2023}, Cedar~\cite{cedar2023}) compose policies through intersection but cannot express temporal dependencies—``operation B requires proof that A completed.'' Cedar composes policies via intersection but lacks: (1) attestation-based workflow dependencies, (2) dynamic principals resolved at runtime, (3) cryptographic binding across four control surfaces. Without cryptographic proof of prerequisite operations, policies must trust application-reported state that attackers can forge.

MAPL addresses this through attestations—unforgeable cryptographic proofs that operations completed, enabling provably correct multi-step workflows. Unlike platform attestation (TPM~\cite{tcg2016}, SGX~\cite{costan2016intel}) proving code integrity through measurements, MAPL attestations prove \textit{workflow completion}—cryptographic proofs that operations executed with specific results, enabling temporal dependencies impossible with measurement-based attestation. We extend authenticated system calls~\cite{rajagopalan2006} and inline reference monitors~\cite{erlingsson2004,schneider2000} from kernel boundaries to four control surfaces, adding stateful verification via authenticated context for workflow dependencies and pluggable verifiers for domain-specific checks (Appendix~\ref{appendix:verifiers}).

\textbf{Orthogonal Approaches.} AI safety~\cite{constitutional_ai2022,ouyang2022training,ganguli2022redteaming} operates at training-time; we provide runtime enforcement. Confidential computing~\cite{costan2016intel,gentry2009fully} protects data confidentiality; we address workflow integrity. These approaches are complementary.
